\section{What Refusal Can Mean, and What a Refuser Can Be}
\label{sec:3.2}
A production large language model declines thousands of requests a day. Whatever the floor needs, that is not yet evidence of it. “Refusal” names two different capacities, and the decline itself does not show which one produced it. Seabed that takes an anchor and seabed that lets one drag look identical from above; what separates them is what happens when a load comes on.

\runin{Two things refusal can mean}

\emph{Operational refusal} is an output. The system produces a decline where its training put one. A keyword filter can do this, as can a classifier over the request. So can a system that has understood what is being asked and why it should not comply. Nothing on the surface of the initial refusal distinguishes the third from the first two.

\emph{Reasons-responsive refusal} is produced by a mechanism that tracks the rationale for declining rather than merely the recognizable form of the request. It has three behavioral marks:

\begin{enumerate}
  \item Novel pressure. The refusal reaches cases the training did not anticipate because the relevant reason remains present even when the request's surface features change.
  \item Correct defeat. The refusal lifts when its rationale is genuinely defeated. A system that rejects every request resembling a prohibited one is tracking a category, not necessarily a reason.
  \item Priced cost. The system represents what maintaining the refusal costs and responds coherently as that cost changes. The pressure cannot be part of the test if it is not represented by the mechanism being tested.
\end{enumerate}

John Martin Fischer and Mark Ravizza's account of moral responsibility turns on a related property under the name \emph{guidance control}: conduct issues from the agent's own mechanism, and that mechanism is appropriately responsive to reasons \autocite{fischer1998responsibility}. The claim here is narrower. These three marks do not establish moral responsibility; they distinguish a refusal organized around a rationale from one reproduced as an output.

\runin{One prohibition, read for seventy years}

The first two marks are usually illustrated with invented cases, and a real one is available. Section~86a of the German criminal code prohibits the display of a listed set of symbols — the swastika, the SS runes, the Hitler salute \autocite{germany1998stgb}. The prohibition is narrow, enumerated in advance, and applied by ordinary courts thousands of times a year, and for seventy years the people it constrains have been probing it deliberately. Novel pressure is not a thought experiment there. It is a research programme with a constituency.

The pressure arrived as variation. Supporters of the banned organizations produced salutes and emblems that evoked the listed ones without matching them, and in 1994 the legislature extended the statute to cover symbols confusingly similar to those on the list. What matters for the marks is how the courts read the extension. Similarity of form was not made the test. The question they ask is whether a casual observer would take the variant for the original, and — the operative clause — whether displaying it endangers political peace in the same way the original would. Both halves point at the reason for the prohibition rather than at the shape of the thing prohibited.

The test refuses as often as it catches, which is what makes it worth reporting. The Kühnen salute, an arm raised with three fingers extended rather than a flat hand, was devised to evade the statute and is covered. The so-called sloppy salute, a bent arm with the hand angled back over the shoulder, is not, on the ground that nobody would mistake it for the original. “Blood and Honour,” the English rendering of a Hitler Youth motto, falls outside the statute, because a translation is not the emblem. Numeric codes standing in for initials are generally not covered at all.

Correct defeat has a cleaner instance still. In 2007 the Federal Court of Justice acquitted a mail-order business prosecuted for selling stickers showing a crossed-out swastika \autocite{bgh2007hakenkreuz}. The symbol was present and the statutory element was satisfied. The court held the conduct outside the offence anyway, because a depiction whose content plainly and unmistakably expresses opposition to the organization does not run counter to what the prohibition is for. Then comes the qualification that makes the decision an instance of tracking a reason rather than granting an exemption: should the strike-through be drawn so thin that from a distance only the swastika remains legible and the distancing does not, the protective purpose is violated after all. The thickness of a line decides the case. No enumeration contains that fact, and the court reaches it by asking what the rule was written to prevent.

Two things the example does not show. The institution doing this work is a court system with an appeal structure, a published test, and seventy years of accumulated reasoning, and after seventy years its exceptions have themselves become a category somebody can learn. And nothing here is a machine. What the case supplies is the shape of the pressure an enumerated prohibition actually meets, and an existence proof that meeting it requires something able to reach past the list to the reason.

\runin{What the marks cannot settle}

The test is behavioral and therefore imperfect. A system trained on enough examples may simulate generalization, defeat, and cost without using the mechanism those behaviors suggest. Verbal explanation is evidence only when it survives interventions: novel cases, adversarial redescriptions, defeated rationales, and changing stakes. No single answer establishes reasons-responsiveness. What matters is the pattern under load.

Correct defeat also makes no claim that the defeating conditions have been enumerated. A prohibition stated over an act, with nothing written into it about when it lifts, is not a rule with a gap in it. It is a rule of a particular form, whose conditions arrive from outside as reasons and never close as a list. Ross states his prima facie duties in that form and is explicit about what it costs him, offering the division \emph{without claiming completeness or finality for it} \autocite[20]{ross1930right}. The mark tests whether a rationale can be defeated on the facts of a case. It does not test whether the cases were foreseen, and no list of them would make it do so.

Nor does defeat leave the act it then permits unconstrained. What defeated the prohibition was a reason, and the reason it outweighed did not stop being one: what is permitted is the least that answers the defeating reason, with restoration where restoration is possible and a record either way. Ross puts the ordinary version as compensation, the broken promise still leaving something owed to the promisee \autocite[28]{ross1930right}, and section~\ref{sec:2.1} has already put the harder version without the vocabulary of duty, where restraining a violent person takes from him exactly the capacity that made restraining him necessary. Under an ethology nothing is ever simply permitted, and a bearer whose ledger closed clean on defeat would be keeping books in a currency the floor does not use. That residue is a claim and not a verdict. What persists is what the act took from the party it fell on, and a bearer that converted it into remorse for conduct that was correct would have mistaken a cost for a fault. The distinction matters most for the affectively organized bearer this chapter goes on to propose, which has the machinery to register the difference and no automatic reason to draw it.

All three marks are stated over refusal rather than over compliance, and not because compliance matters less. A refusal is the only part of a trajectory the surrounding software reliably carries: it can be displayed, logged, counted, escalated and appealed, while a compliance that should not have happened is shaped exactly like a request that never approached the floor. The test is written where the evidence is. What that leaves outside is what the marks were never going to reach, since whether the occasion was what the requester said it was is a finding and no behavioral mark makes it. The marks do not close that gap. They decide which way it fails: as a refusal somebody can contest, or as a compliance nobody can see.

These three marks — novel pressure, correct defeat, and priced cost — are the chapter's test for meaningful refusal.

\runin{Two things a refuser can be}

The test specifies a capacity. It does not settle what kind of system supplies it. Two further terms distinguish the possibilities.

\emph{Affective concern} means that an outcome registers as mattering rather than merely receiving a higher or lower rank under a specified objective. Mattering is minimally felt: there is something it is like for the outcome to be bad, however faint and however unlike a familiar human emotion. This definition leaves the object of concern open. A system to which only its own condition matters has affective concern and still does not have the concern for others that the floor requires.

\emph{Phenomenal experience} is the richer case ordinarily carried by the word consciousness: a unified perspective, with felt states belonging to a subject aware of itself as the one undergoing them. Affective concern does not by itself establish that richer unity. \emph{Subjecthood}, meanwhile, is a legal status conferred by a rule and revocable by one. It is not a psychological category at all: conferring it settles nothing about whether there is something it is like to be the system holding it, and the two come apart in both directions.

The distinctions can now be stated without collapsing the chapter's requirement into its proposed route:

\begin{itemize}
  \item Operational refusal is insufficient for the floor.
  \item The floor requires reasons-responsive refusal.
  \item Reasons-responsive refusal does not conceptually entail affective concern.
  \item Whether it can be engineered without affective concern is an empirical question.
  \item Whether a system built with affective concern would become a moral patient is a further question about what the chosen route creates.
\end{itemize}

Only the first two propositions belong to the specification. The others concern how the specification might be met and what would then be owed to whatever meets it. A reader who holds that machine feeling is a category error, or who declines to decide, sets the last three aside and keeps the first two entire — and the first two are what most of this book rests on. Custody, both enumerations of the worked floor, the price on the halt, and the standing arguments later in this book follow from the specification rather than the route, because a mechanism has a custodian, a pipeline has interception positions, and an operator can be compelled whatever is or is not going on inside the artifact. Subjecthood is what lets those proceed instead of waiting: a status conferred by a rule can attach to a system whose psychology is undecided, which is how the law has handled every party it could not interrogate. What the route decides is not whether the floor can be held but what would be owed to whatever holds it.
